% introduction.tex

\section{Introduction}
\label{sec:introduction}

In standard costly signaling theory, separation survives because low types cannot afford to imitate high types. Education, credentials, public effort, and other ambitious signals are informative because the bad apple cannot climb up. This paper studies the opposite force. A high type may be able to send the ambitious signal, and may currently benefit from doing so, yet rationally abandon it because the same signal exposes her to future scrutiny. The good apple climbs down.

The mechanism is not ordinary risk aversion, nor a simple increase in signal cost. The ambitious signal does two things at once. First, it conveys ability. Second, it creates a public claim about the agent's self-model: by choosing the ambitious action, the agent appears to believe that she can execute it. If later scrutiny becomes more severe, the signal is not merely evaluated as performance evidence. It is also repriced as evidence about the agent's self-understanding. The observer judges not only what the agent is, but what the agent seemed to believe about what she was.

This paper formalizes that logic in a dynamic psychological signaling game. There are two signals: an ambitious signal, \(s_A\), and a pooling-safe signal, \(s_P\). There are two types, \(\theta_H\) and \(\theta_L\). The judgment environment follows an exogenous Markov process \(j_t\in\{L,H\}\), where \(L\) denotes low scrutiny and \(H\) denotes high scrutiny. The transition probability
\[
q=\Pr(j_{t+1}=H\mid j_t=L)
\]
is observed before the agent chooses a signal. Sending \(s_A\) creates an exposure state. If high scrutiny arrives in the future, this exposure is socially repriced as a second-order claim about the accuracy of the agent's earlier self-model.

The central result is a top-down pooling reversal. In the low-scrutiny state, a high type would choose the ambitious signal under the standard static logic. However, if the probability \(q\) of entering high scrutiny becomes large enough, the high type internalizes the future epistemic penalty and abandons the ambitious signal. The separating profile fails because the high type's own incentive constraint breaks. Pooling does not arise because low types flood the high signal. It arises because high types exit the high signal.

The force behind the result is a scrutiny-type complementarity generated by adverse repricing. The primitive object is not signal precision and not a history-induced posterior. Each type has a primitive self-model-accuracy level
\[
p_\theta=\Pr(m=1\mid \theta),
\]
where \(m=1\) denotes that the sender's self-model is accurate. Higher types have higher self-accuracy levels: \(p_H>p_L\). High scrutiny is not a truth audit that mechanically confirms accurate high types. It is a judgment-intensity state in which exposed self-presentations can be adversely reinterpreted. When that adverse reinterpretation arrives, the relevant loss is the Bayesian-surprise distance between the type's carried self-accuracy level and a lower adverse benchmark \(\bar p\). Thus higher types are not more exposed because the receiver happens to label them high on the equilibrium path. They are more exposed because their true type carries more self-model reputation-at-risk.

This distinction matters for equilibrium selection. In standard signaling models, off-path ambitious signals are credible because high types are more willing to send them. Forward-induction refinements such as the intuitive criterion and D1 restrict off-path beliefs toward the type with the stronger incentive to deviate. Here, future scrutiny reverses that willingness ranking. When \(q\) is high, the high type is more reluctant to send \(s_A\) because her future reputation-at-risk is larger. A deviation to the ambitious signal is then attributed not to the confident high type, but to the low or poorly calibrated type who fails to internalize future exposure. The ambitious signal loses its high-type interpretation.

The paper contributes to three literatures. First, it extends costly signaling theory by identifying a failure mode that is the mirror image of the standard Spence logic. In \citet{spence1973job}, separation is supported because low types cannot profitably mimic high types. In the present model, separation fails because high types no longer wish to occupy the separating position. The D1 logic of \citet{cho1987signaling}, \citet{banks1987equilibrium}, and \citet{cho1990strategic} is not rejected; it is reversed by anticipatory exposure.

Second, the paper contributes to psychological game theory. Dynamic psychological games allow payoffs to depend on higher-order beliefs, including beliefs about how others interpret one's intentions, self-image, or plans \citep{battigalli2009dynamic}. The present model uses this language to study a specific object: second-order self-model inference. The social cost of a signal depends on what observers infer the agent believed about herself when she chose it.

Third, the paper connects to theories of self-confidence and ego utility. \citet{benabou2002self} emphasize that self-confidence can have consumption, signaling, and motivation value. \citet{koszegi2006ego} studies how ego utility can distort task choice toward ambitious actions. The present paper studies the dynamic downside of that same self-presentational structure. A signal that is valuable under low scrutiny can become a liability when future judgment intensity rises.

The model also yields an observable implication. Near the critical transition probability, the quality of the ambitious signal need not simply mean-revert. In a continuous-type extension, ordinary winners withdraw first because they have enough reputation-at-risk to be exposed but not enough surplus to withstand high scrutiny. The remaining senders of \(s_A\) can be polarized: extremely strong types continue, while poorly calibrated low types also remain. The signal is hollowed out from the middle.

The rest of the paper proceeds as follows. Section~\ref{sec:model} introduces the primitives, the judgment-intensity process, and the exposure technology. Section~\ref{sec:surprise} defines type-indexed self-accuracy levels and derives scrutiny-type complementarity from Bayesian-surprise costs under adverse repricing. Section~\ref{sec:dynamic} characterizes the anticipatory exposure value, the top-down threshold, and the D1 flip. Section~\ref{sec:polarization} studies the continuous-type extension and derives polarized degradation of ambitious signals.
